
\section{Related Work}
\subsection{Advertising Ecosystem Auditing}
The rise of transparency requirements for digital advertising has led to efforts in measuring how effectively companies are in following this goal \cite{metaAudit}. Auditing becomes the object of study itself, with policies compared between platforms, regulations and public opinion \cite{politicalComplexity}. A study focusing on Facebook has shown how algorithmic advert delivery optimization can lead to discriminatory content distribution \cite{adDeliveryDiscrimination}.

The abuse of otherwise legitimate distribution systems for malicious purposes has been a research topic in several studies. In \cite{adFraudRevenue}, authors investigate how identifier pooling can be used to redirect revenue to fraudulent sources. Through the analysis of four different traffic sources, a study by \cite{urlShortAbuse_Cloaking_WWW} investigates abusive advertising distribution, revealing that cloaking-aware scraping leads to the discovery of 81\% more malicious URLs than singular crawling strategies, also highlighting the relevance of URL shorteners in fraudulent advertising. 

Meta's Ad Library has been extensively studied in the context of political advertising ~\cite{adpolit1, adpolit2, adpolit3, adpolit4}, with the ecosystem also being investigated externally \cite{politAuditMetaWWW, transparencyLimitationAdLib} from the library to further transparency efforts. It has also been studied in the context of Illicit Betting Ecosystems \cite{adLibBet}, Health-related topics, such as ``alternative'' medicine promotion \cite{healthMetaAd} and COVID-19 informational weaponization \cite{adLibHealth}. One security-centered study points to design and implementation flaws that harm the Library's transparency, revealing the presence of undeclared, coordinated advertising campaigns associated with inauthentic content \cite{adlibSecAudit}. 

\subsection{Cloaking}
With over 20 years both as a conceptual research topic and as a set of strategies for online deception, the term's usage has evolved with the different, emerging strategies that enable cloaking in different contexts. 

\textbf{Redirection Cloaking}: . Research by \cite{redirectionFormats} provides a general understanding of redirection mechanisms used in cloaking, outlining the same practices we empirically observe.

\textbf{Url Masking}: A study by \cite{adUrlRisksRelated} analyzed the security risks of ad-based URL shortening services. 

\textbf{Creative Optimization}: The exploitation of dynamic content delivery systems avoid moderation has been a phenomenon reported in a wide range of platforms, being confirmed as an effect bypass for contextual ads in Youtube \cite{youtubeContextCloak}



Currently, the closest a study has gotten to analyzing cloaking patterns in an online ecossystem is \cite{industrialCloaking}, which does not mention the specific platform.



\subsection{Research Gap}
Despite plentiful related analyses, research has frequently focused on context-specific definitions \cite{urlShortAbuse_Cloaking_WWW}, leading to rich, highly diverse reporting on actively employed cloaking techniques. However, such diversity has also led to similar phenomena not being considered under an unifying framewwork.

With, to the best of our knowledge, no previous studies on advert-embedded URLs on Meta's advertising ecossystem, our study aims to fill in:
\begin{enumerate}
    \item A conceptual gap on the definition of cloaking through an analysis of its occourence on in one of the most relevant advertising ecossystems in the world.
    \item A gap in security auditing literature where no large-scale analyses have been condutcted in relation to URL embeds for adverts in Metas's Ad Library.
\end{enumerate}



%
\begin{table}[b]
\resizebox{\columnwidth}{!}{%
\begin{tabular}{@{}lll@{}}
\toprule
\textbf{Paper}                                                                                                                                 & \textbf{Venue} & \textbf{Reference for}                                                                                                                                                                                                  \\ \midrule
\begin{tabular}[c]{@{}l@{}}Welcome to the Dark Side:\\ Revenue Flows of Fraud in the Online Ad Ecosystem\end{tabular}                          & www 25         & \textbf{\begin{tabular}[c]{@{}l@{}}transparency / ad delivery mechanism\\ \\ remoção retroativa\end{tabular}}                                                                                                           \\
\begin{tabular}[c]{@{}l@{}}Altering Words to Evade Perceived Moderation:\\ Decoding Algospeak\end{tabular}                                     & icwsm 26       & \begin{tabular}[c]{@{}l@{}}evolução do termo Cloaking\\ (não usa o termo diretamente mas é análogo)\end{tabular}                                                                                                        \\
Discrimination through Optimization                                                                                                            & cscw 19        & \textbf{transparency / ad delivery mechanism}                                                                                                                                                                           \\
\begin{tabular}[c]{@{}l@{}}50 Shades of Deceptive Patterns: A Unified Taxonomy,\\ Multimodal Detection, and Security Implications\end{tabular} & www 25         & \begin{tabular}[c]{@{}l@{}}taxonomy\\ \\ similar to our paper on a different context\end{tabular}                                                                                                                       \\
\begin{tabular}[c]{@{}l@{}}Where are you taking me? \\ Understanding Abusive Traffic Distribution Systems\\ (ODIN)\end{tabular}                & www 21         & \begin{tabular}[c]{@{}l@{}}simultaneous cloaking \\ (typosquat, piracy, ad-shorteners, scam landing)\\ \\ url shortening abuse\end{tabular}                                                                             \\
Short Path to Phishing                                                                                                                         & eCrime 25      & \begin{tabular}[c]{@{}l@{}}url shortening abuse\\ \\ lexical features\\ \\ redirection mechanisms\end{tabular}                                                                                                          \\
A Security Analysis of the Facebook Ad Library                                                                                                 & IEEE S\&P 20   & \textbf{\begin{tabular}[c]{@{}l@{}}transparency / ad delivery mechanism\\ \\ inauthentic coordination\end{tabular}}                                                                                                     \\
Needle in a Haystack                                                                                                                           & IMC 2018       & \begin{tabular}[c]{@{}l@{}}224M records -\textbackslash{}textgreater 660k squatting -\textbackslash{}textgreater 1.2k phish\\ (rarity of phishing domains)\\ \\ \\ claims ocr + vision defeats obfuscation\end{tabular} \\ \bottomrule
\end{tabular}%
}
\end{table}

